\documentclass{article}
\usepackage{ctex}
\usepackage{amsmath}
\usepackage{amssymb}
\usepackage{cite}

\newtheorem{example}{例}             % 整体编号
\newtheorem{theorem}{定理}[section]  % 按 section 编号
\newtheorem{definition}{定义}
\newtheorem{axiom}{公理}
\newtheorem{property}{性质}
\newtheorem{proposition}{命题}
\newtheorem{lemma}{引理}
\newtheorem{corollary}{推论}
\newtheorem{remark}{注解}
\newtheorem{condition}{条件}
\newtheorem{conclusion}{结论}
\newtheorem{assumption}{假设}


\title{Mandelbrot 集及其算法实现}
\author{杨子宸}

\begin{document}

\maketitle

\begin{abstract}
本文介绍了 Mandelbrot 集的定义和算法实现, 通过简单的迭代方法, 我们可以在计算机上生成 Mandelbrot 集的图像.
\end{abstract}

\section{引言}
Mandelbrot 集合作为分形几何中最著名的研究对象之一, 由 Benoit Mandelbrot 首次提出\cite{mandelbrot1982fractal}. 这个定义看似简单的集合, 在复平面上却展现出无限精细的复杂结构和自相似特性. 本文将介绍 Mandelbrot 集合的数学原理, 以及如何在计算机上实现其图像生成.

\section{Mandelbrot集的数学原理}

\begin{definition}
对于非线性迭代序列
\[z_{n + 1} = z_n^2 + c, z_0 = 0,\]
其中\(c \in \mathbb{C}\), \(z_n \in \mathbb{C}\),定义 Mandelbrot 集为当\(n \to \infty\)时所有使得序列\(z_n\)保持有限 (特别地, \(|z_n| \leq 2\)) 的\(c\)的集合, 即
\[\mathcal{M} = \{c \in \mathbb{C}: |z_n| \leq 2\;(\forall n\in\mathbb{N})\}.\]
\end{definition}

\begin{property}
\label{prop::bounded}
设\(c \in \mathbb{C}\), \(z_n(c) := \left\{\begin{aligned}
&z_0(c) = 0,\\
&z_{n + 1}(c) = z_n(c)^2 + c, \quad n \geq 0.
\end{aligned}\right.\) 若\(c \in \mathcal{M}\), 则\(\forall n \in \mathbb{N}\):
\[|z_n(c)| \leq 2.\]
\end{property}

证明见参考文献\cite{branner1989mandelbrot}. 这个性质是 Mandelbrot 集生成算法的原理. 

\section{Mandelbrot 集的算法实现}
绘制 Mandelbrot 集合一般采用 ``逃逸时间算法'', 步骤如下:

\begin{enumerate}
\item 确定复平面上的范围和分辨率.
\item 对每个复数\(c\), 从\(z = 0\)开始迭代公式\(z_{n + 1} = z_n^2 + c\). 
\item 由性质\ref{prop::bounded}, 如果在最大迭代次数内\(|z| > 2\), 则认为\(c \notin \mathcal{M}\). 
\item 根据逃逸速度为不属于集合的点赋予颜色. 
\end{enumerate}

\section{结论}
Mandelbrot 集被称为 ``上帝的指纹'' , 是科学与艺术的美妙结合, 它的形状不仅十分美观, 也蕴含着深刻的数学内涵.

\bibliography{mandelbrot_ref.bib}
\bibliographystyle{plain}
\end{document}
